% ===== PRÉAMBULE CANONIQUE — à coller VERBATIM en tête de CHAQUE .tex =====
\documentclass[11pt,a4paper]{article}
\usepackage[utf8]{inputenc}\usepackage[T1]{fontenc}\usepackage{lmodern}
\usepackage[french]{babel}
\usepackage{amsmath,amssymb,mathtools}\usepackage{icomma}
\usepackage{siunitx}\sisetup{locale=FR}  % \qty{}{}, \si{}, \ang{}
\usepackage{xcolor}\usepackage{colortbl}
\usepackage{tikz}\usepackage{pgfplots}\pgfplotsset{compat=1.18}
\usepackage[siunitx,europeanresistors]{circuitikz} % circuits (élec.)
\usepackage[most]{tcolorbox}\usepackage{enumitem}\usepackage{array,booktabs}
\usepackage[a4paper,margin=2.2cm]{geometry}
\usetikzlibrary{arrows.meta,calc,angles,quotes,positioning,patterns,%
  backgrounds,shapes.geometric,decorations.pathmorphing,decorations.markings,intersections}
\definecolor{primary}{RGB}{222,49,99}\definecolor{primarydark}{RGB}{150,22,55}
\definecolor{defcol}{RGB}{0,114,178}\definecolor{methcol}{RGB}{52,92,148}
\definecolor{excol}{RGB}{0,135,90}\definecolor{attncol}{RGB}{200,40,55}
\tcbset{boxbase/.style={enhanced,breakable,boxrule=0.9pt,arc=2.5pt,
  left=3mm,right=3mm,top=2.3mm,bottom=2.3mm,fonttitle=\bfseries,coltitle=white}}
% --- boîtes TUTORIEL (noms EXACTS, ne pas renommer) ---
\newtcolorbox{cle}[1][]{boxbase,colback=defcol!6,colframe=defcol,
  title={\textsc{À retenir}\ifx\relax#1\relax\relax\else\textmd{ : \itshape #1}\fi}}
\newtcolorbox{methode}[1][]{boxbase,colback=methcol!7,colframe=methcol,sharp corners,
  title={\textsc{Méthode}\ifx\relax#1\relax\relax\else\textmd{ --- \itshape #1}\fi}}
\newtcolorbox{exemplebox}[1][]{boxbase,colback=excol!5,colframe=excol,
  title={\textsc{Exemple}\ifx\relax#1\relax\relax\else\textmd{ : \itshape #1}\fi}}
\newtcolorbox{piege}[1][]{boxbase,colback=attncol!6,colframe=attncol,
  title={\textsc{Piège}\ifx\relax#1\relax\relax\else\textmd{ : \itshape #1}\fi}}
\newtcolorbox{remarque}[1][]{boxbase,colback=black!4,colframe=black!45,
  title={\textsc{Remarque}\ifx\relax#1\relax\relax\else\textmd{ : \itshape #1}\fi}}
% --- boîtes EXERCICE / CORRIGÉ (noms EXACTS, auto-comptées) ---
\newtcolorbox[auto counter]{exo}[1][]{boxbase,colback=primary!4,colframe=primary,
  title={\textsc{Exercice}~\thetcbcounter\ifx\relax#1\relax\relax\else\textmd{ --- \itshape #1}\fi}}
\newtcolorbox[auto counter]{corrige}[1][]{boxbase,colback=excol!4,colframe=excol,
  title={\textsc{Corrigé}~\thetcbcounter\ifx\relax#1\relax\relax\else\textmd{ --- \itshape #1}\fi}}
\newcommand{\vv}[1]{\vec{#1}}\newcommand{\ds}{\displaystyle}
\newcommand{\R}{\mathbb{R}}\newcommand{\N}{\mathbb{N}}\newcommand{\Z}{\mathbb{Z}}
% (un \usepackage{fancyhdr}+en-tête est possible ; il est ignoré côté web)
% =========================================================================
\begin{document}
\sisetup{per-mode=symbol}

\begin{exo}[champ d'une charge ponctuelle]
Une charge $Q = \SI{+8}{\nano\coulomb}$ est isolée dans l'air.
\begin{enumerate}[label=\alph*)]
  \item Calcule l'intensité du champ électrique qu'elle crée à $r = \SI{20}{\centi\metre}$.
  \item Dans quel sens pointe $\vv{E}$ en ce point (vers la charge ou à l'opposé) ?
\end{enumerate}
\end{exo}

\begin{exo}[de la force au champ]
Une charge test $q = \SI{+4}{\nano\coulomb}$ placée en un point $M$ subit une force électrique
d'intensité $F = \SI{8e-6}{\newton}$.
\begin{enumerate}[label=\alph*)]
  \item Calcule l'intensité du champ électrique $E$ en $M$.
  \item Si l'on remplace la charge test par $q' = \SI{+8}{\nano\coulomb}$, que vaut le champ en $M$ ?
\end{enumerate}
\end{exo}

\begin{exo}[la force dans un champ]
En un point règne un champ uniforme $E = \SI{500}{\newton\per\coulomb}$, dirigé vers la droite.
\begin{enumerate}[label=\alph*)]
  \item Quelle force (intensité et sens) subit une charge $q = \SI{+2}{\micro\coulomb}$ ?
  \item Même question pour une charge $q = \SI{-2}{\micro\coulomb}$.
\end{enumerate}
\end{exo}

\begin{exo}[orienter le champ]
On considère une charge positive et une charge négative.
\begin{center}
\begin{tikzpicture}[scale=1,>=Stealth,font=\footnotesize,
    ch/.style={circle,minimum size=8mm,inner sep=0pt,text=white,font=\bfseries\small}]
  \node[ch,fill=attncol] (P) at (-3,0) {$+$};
  \node[circle,fill=black,inner sep=1.2pt,label=below:{$M_1$}] at (-1.4,0) {};
  \node[ch,fill=defcol] (N) at (2.6,0) {$-$};
  \node[circle,fill=black,inner sep=1.2pt,label=below:{$M_2$}] at (1.0,0) {};
\end{tikzpicture}
\end{center}
\begin{enumerate}[label=\alph*)]
  \item Dessine (ou décris) le vecteur $\vv{E}$ créé par la charge $+$ au point $M_1$.
  \item Dessine (ou décris) le vecteur $\vv{E}$ créé par la charge $-$ au point $M_2$.
\end{enumerate}
\end{exo}

\begin{exo}[champ nul au milieu]
Deux charges \emph{identiques} $q_A = q_B = +q$ sont placées en $A$ et $B$, distants de $2a$.
On note $C$ le milieu de $[AB]$.
\begin{center}
\begin{tikzpicture}[scale=1,>=Stealth,font=\footnotesize,
    ch/.style={circle,fill=attncol,minimum size=8mm,inner sep=0pt,text=white,font=\bfseries\small}]
  \node[ch] (A) at (-2.4,0) {$+q$}; \node[below,font=\scriptsize] at (-2.4,-0.45) {$A$};
  \node[ch] (B) at (2.4,0) {$+q$}; \node[below,font=\scriptsize] at (2.4,-0.45) {$B$};
  \node[circle,fill=black,inner sep=1.2pt] at (0,0) {}; \node[below,font=\scriptsize] at (0,-0.4) {$C$};
\end{tikzpicture}
\end{center}
\begin{enumerate}[label=\alph*)]
  \item Compare les intensités des champs $\vv{E}_A$ et $\vv{E}_B$ créés en $C$.
  \item En déduire le champ total en $C$. Justifie par les sens.
\end{enumerate}
\end{exo}

\end{document}
